\documentclass[11pt]{article}
% The preamble every project's paper shares: % The preamble every project's paper shares: % The preamble every project's paper shares: \input{../../preamble} from
% projects/<project>/paper/. Boards come from chessboard and xskak, plots from
% pgfplots; every number a paper prints comes from its generated.tex.
\usepackage[LSBC4,T1]{fontenc}
\usepackage[utf8]{inputenc}
\usepackage{lmodern}
\usepackage[margin=1in]{geometry}
\usepackage{microtype}
\usepackage{amsmath}
\usepackage{booktabs}
\usepackage{tabularx}
\usepackage{array}
\usepackage{xcolor}
\usepackage[obeyspaces,spaces]{url}
\usepackage{xskak}
\usepackage{pgfplots}
\pgfplotsset{compat=1.18}
\usepackage{caption}
\usepackage{subcaption}
\usepackage{enumitem}
\usepackage[section]{placeins}
\usepackage[hidelinks]{hyperref}

\setchessboard{boardfontencoding=LSBC4, boardfontfamily=skaknew, blackfieldcolor=black!22,
  whitefieldcolor=white, setfontcolors, showmover=false, label=true, labelfontsize=6pt,
  boardfontsize=14pt}

% A move or line in UCI, and a position in FEN, both breakable at a slash.
\newcommand{\uci}[1]{\texttt{#1}}
\newcommand{\fen}[1]{{\footnotesize\path{#1}}}
\newcolumntype{L}{>{\raggedright\arraybackslash}X}

\setlength{\parskip}{0.4em}
\setlength{\parindent}{0pt}
 from
% projects/<project>/paper/. Boards come from chessboard and xskak, plots from
% pgfplots; every number a paper prints comes from its generated.tex.
\usepackage[LSBC4,T1]{fontenc}
\usepackage[utf8]{inputenc}
\usepackage{lmodern}
\usepackage[margin=1in]{geometry}
\usepackage{microtype}
\usepackage{amsmath}
\usepackage{booktabs}
\usepackage{tabularx}
\usepackage{array}
\usepackage{xcolor}
\usepackage[obeyspaces,spaces]{url}
\usepackage{xskak}
\usepackage{pgfplots}
\pgfplotsset{compat=1.18}
\usepackage{caption}
\usepackage{subcaption}
\usepackage{enumitem}
\usepackage[section]{placeins}
\usepackage[hidelinks]{hyperref}

\setchessboard{boardfontencoding=LSBC4, boardfontfamily=skaknew, blackfieldcolor=black!22,
  whitefieldcolor=white, setfontcolors, showmover=false, label=true, labelfontsize=6pt,
  boardfontsize=14pt}

% A move or line in UCI, and a position in FEN, both breakable at a slash.
\newcommand{\uci}[1]{\texttt{#1}}
\newcommand{\fen}[1]{{\footnotesize\path{#1}}}
\newcolumntype{L}{>{\raggedright\arraybackslash}X}

\setlength{\parskip}{0.4em}
\setlength{\parindent}{0pt}
 from
% projects/<project>/paper/. Boards come from chessboard and xskak, plots from
% pgfplots; every number a paper prints comes from its generated.tex.
\usepackage[LSBC4,T1]{fontenc}
\usepackage[utf8]{inputenc}
\usepackage{lmodern}
\usepackage[margin=1in]{geometry}
\usepackage{microtype}
\usepackage{amsmath}
\usepackage{booktabs}
\usepackage{tabularx}
\usepackage{array}
\usepackage{xcolor}
\usepackage[obeyspaces,spaces]{url}
\usepackage{xskak}
\usepackage{pgfplots}
\pgfplotsset{compat=1.18}
\usepackage{caption}
\usepackage{subcaption}
\usepackage{enumitem}
\usepackage[section]{placeins}
\usepackage[hidelinks]{hyperref}

\setchessboard{boardfontencoding=LSBC4, boardfontfamily=skaknew, blackfieldcolor=black!22,
  whitefieldcolor=white, setfontcolors, showmover=false, label=true, labelfontsize=6pt,
  boardfontsize=14pt}

% A move or line in UCI, and a position in FEN, both breakable at a slash.
\newcommand{\uci}[1]{\texttt{#1}}
\newcommand{\fen}[1]{{\footnotesize\path{#1}}}
\newcolumntype{L}{>{\raggedright\arraybackslash}X}

\setlength{\parskip}{0.4em}
\setlength{\parindent}{0pt}

% Written by projects/drill-decks/paper/generate.mjs; do not edit by hand.

\newcommand{\AppCommit}{d8b771238c}
\newcommand{\Items}{1,400}
\newcommand{\Scored}{1,336}
\newcommand{\Deviations}{64}
\newcommand{\Defensive}{489}
\newcommand{\Attacking}{286}
\newcommand{\Mixed}{561}
\newcommand{\Both}{84}
\newcommand{\Neither}{477}
\newcommand{\NeitherShare}{36}
\newcommand{\DefensiveShare}{37}
\newcommand{\AttackingShare}{21}
\newcommand{\PhaseItems}{781}
\newcommand{\PhaseShare}{56}
\newcommand{\MotifItems}{150}
\newcommand{\MotifShare}{11}
\newcommand{\DefItems}{405}
\newcommand{\DefShare}{29}
\newcommand{\FeedsSkill}{742}
\newcommand{\FeedsSkillShare}{56}
\newcommand{\ThemeTable}{%
\begin{tabular}{llrrrrl}
\toprule
word & skill it feeds & items & defensive & attacking & mixed & fits \\
\midrule
\texttt{defensiveMove} & \texttt{defensiveMove} & 405 & 195 & 56 & 154 & 195 of 405 (48\%) \\
\texttt{hangingPiece} & \texttt{hangingPiece} & 109 & 13 & 54 & 42 & 70 of 109 (64\%) \\
\texttt{fork} & \texttt{fork} & 20 & 1 & 15 & 4 & 16 of 20 (80\%) \\
\texttt{mateIn1} & \texttt{mateIn1} & 20 & 0 & 17 & 3 & 20 of 20 (100\%) \\
\texttt{backRankMate} & \texttt{backRankMate} & 1 & 0 & 0 & 1 & 0 of 1 (0\%) \\
\texttt{middlegame} & -- & 406 & 139 & 81 & 186 & not scored \\
\texttt{endgame} & \texttt{conv.technique} & 187 & 68 & 38 & 81 & not scored \\
\texttt{opening} & -- & 188 & 73 & 25 & 90 & not scored \\
\bottomrule
\end{tabular}
}
\newcommand{\DefFit}{48}
\newcommand{\DefFitK}{195}
\newcommand{\DefAttacking}{56}
\newcommand{\DefensiveElsewhere}{294}
\newcommand{\DefRecall}{40}
\newcommand{\PhaseAttacking}{144}
\newcommand{\PhaseDefensive}{280}
\newcommand{\AttackingNamed}{30}
\newcommand{\PhaseAttackingNothing}{61}
\newcommand{\EndgameItems}{187}
\newcommand{\EndgameDefensive}{68}
\newcommand{\EndgameCompetency}{\texttt{conv.technique}}
\newcommand{\HangFit}{64}
\newcommand{\HangPawn}{22}
\newcommand{\HangFitPawn}{84}
\newcommand{\HangItems}{109}
\newcommand{\MateFit}{20 of 20}
\newcommand{\ForkFit}{16 of 20}
\newcommand{\Quizzes}{48}
\newcommand{\Margin}{45}
\newcommand{\TwoRight}{17}
\newcommand{\TwoRightShare}{35}
\newcommand{\OneRight}{31}
\newcommand{\Forcing}{30}
\newcommand{\Material}{eleven}
\newcommand{\ForcingOrMaterial}{31}
\newcommand{\ForcingShare}{65}
\newcommand{\PlayedBest}{34}
\newcommand{\PlanQuestions}{seven}
\newcommand{\PlanShare}{15}
\newcommand{\DefBoard}{%
\chessboard[setfen={r4rk1/pp2np1p/5Qp1/1qbp4/6P1/2Bn3P/PP2NP2/1R3K1R b - - 5 17}, boardfontsize=14pt, pgfstyle=straightmove, arrow=to, linewidth=0.07em, color=red!75!black, markmoves={d3-f2}, pgfstyle=straightmove, arrow=to, linewidth=0.07em, color=green!50!black, markmoves={d5-d4}]
}
\newcommand{\DefFen}{\fen{r4rk1/pp2np1p/5Qp1/1qbp4/6P1/2Bn3P/PP2NP2/1R3K1R b - - 5 17}}
\newcommand{\DefPlayed}{17\ldots{}\symknight{}xf2}
\newcommand{\DefLine}{17\ldots{}d4 18.~\symknight{}xd4 \symbishop{}xd4 19.~\symbishop{}xd4 \symknight{}f4+ 20.~\symking{}g1 \symknight{}e2+ 21.~\symking{}h2}
\newcommand{\DefUci}{\uci{d5d4 e2d4 c5d4 c3d4 d3f4 f1g1 f4e2 g1h2}}
\newcommand{\DefTheme}{\texttt{middlegame}}
\newcommand{\DefSkill}{no skill}
\newcommand{\DefDrop}{86}
\newcommand{\DefPunish}{18.~\symqueen{}g7\#}
\newcommand{\WinBoard}{%
\chessboard[setfen={4r3/p7/1p3p2/8/1P4k1/P1P1B1R1/4nPK1/8 b - - 2 38}, boardfontsize=14pt, pgfstyle=straightmove, arrow=to, linewidth=0.07em, color=red!75!black, markmoves={g4-f5}, pgfstyle=straightmove, arrow=to, linewidth=0.07em, color=green!50!black, markmoves={e2-g3}]
}
\newcommand{\WinFen}{\fen{4r3/p7/1p3p2/8/1P4k1/P1P1B1R1/4nPK1/8 b - - 2 38}}
\newcommand{\WinPlayed}{38\ldots{}\symking{}f5}
\newcommand{\WinLine}{38\ldots{}\symknight{}xg3 39.~a4 \symknight{}e2 40.~f3+ \symking{}f5 41.~\symbishop{}d2 \symknight{}d4 42.~\symbishop{}c1}
\newcommand{\WinUci}{\uci{e2g3 a3a4 g3e2 f2f3 g4f5 e3d2 e2d4 d2c1}}
\newcommand{\WinTheme}{\texttt{endgame}}
\newcommand{\WinSkill}{\texttt{conv.technique}}
\newcommand{\WinDrop}{40}
\newcommand{\WinPunish}{39.~\symrook{}f3+ \symking{}e4 40.~\symbishop{}g5 \symrook{}g8 41.~\symrook{}e3+ \symking{}d5}
\newcommand{\DefWinsSmall}{15}
\newcommand{\DefWins}{34}
\newcommand{\QuizBoard}{%
\chessboard[setfen={rn2kb1r/pp1qp1pp/2p2n2/3p1p2/B2P1B2/2NQPb2/PPP2PPP/R4RK1 w kq - 0 11}, boardfontsize=14pt, pgfstyle=straightmove, arrow=to, linewidth=0.07em, color=green!50!black, markmoves={g2-f3}]
}
\newcommand{\QuizFen}{\fen{rn2kb1r/pp1qp1pp/2p2n2/3p1p2/B2P1B2/2NQPb2/PPP2PPP/R4RK1 w kq - 0 11}}
\newcommand{\QuizLine}{11.~gxf3 e6 12.~e4 \symbishop{}e7 13.~exf5 exf5 14.~\symking{}h1 g6 15.~\symrook{}ae1 \symking{}f7}
\newcommand{\QuizUci}{\uci{g2f3 e7e6 e3e4 f8e7 e4f5 e6f5 g1h1 g7g6 a1e1 e8f7}}
\newcommand{\QuizOptions}{11.~gxf3 (0.2), 11.~\symbishop{}xb8 ($-$2.1), 11.~\symrook{}ae1 ($-$2.3)}
\newcommand{\QuizLead}{2.3}
\newcommand{\ThreatTable}{%
\begin{tabular}{lrrrrrrrr}
\toprule
 & & & \multicolumn{3}{c}{threatened} & \multicolumn{3}{c}{the gate} \\
\cmidrule(lr){4-6}\cmidrule(lr){7-9}
 & positions & in check & $n$ & \% & mate & fires \% & precision \% & recall \% \\
\midrule
clear defensive miss & 489 & 45 & 260 & 59 & 42 & 58 & 75 & 75 \\
clear missed win & 286 & 12 & 86 & 31 & 7 & 50 & 51 & 81 \\
both & 84 & 9 & 34 & 45 & 6 & 57 & 63 & 79 \\
no material change & 477 & 21 & 140 & 31 & 11 & 44 & 55 & 78 \\
\midrule
\texttt{defensiveMove} & 405 & 11 & 271 & 69 & 39 & 100 & 69 & 100 \\
\texttt{hangingPiece} & 109 & 9 & 36 & 36 & 4 & 51 & 57 & 81 \\
\texttt{fork} & 20 & 0 & 6 & 30 & 0 & 60 & 42 & 83 \\
\texttt{mateIn1} & 20 & 0 & 12 & 60 & 1 & 80 & 63 & 83 \\
\texttt{backRankMate} & 1 & 0 & 1 & 100 & 0 & 100 & 100 & 100 \\
\texttt{middlegame} & 406 & 21 & 110 & 29 & 11 & 26 & 53 & 49 \\
\texttt{endgame} & 187 & 37 & 45 & 30 & 10 & 21 & 45 & 31 \\
\texttt{opening} & 188 & 9 & 39 & 22 & 1 & 18 & 52 & 44 \\
\midrule
all & 1336 & 87 & 520 & 42 & 66 & 51 & 63 & 77 \\
\bottomrule
\end{tabular}
}
\newcommand{\Budget}{600,000}
\newcommand{\Searches}{2,585}
\newcommand{\InCheck}{87}
\newcommand{\Searched}{1,249}
\newcommand{\Threatened}{520}
\newcommand{\ThreatShare}{42}
\newcommand{\ThreatMate}{66}
\newcommand{\ThreatMaterial}{454}
\newcommand{\NeitherSearched}{456}
\newcommand{\NeitherThreat}{140}
\newcommand{\NeitherThreatShare}{31}
\newcommand{\DefSearched}{444}
\newcommand{\DefThreat}{260}
\newcommand{\DefThreatShare}{59}
\newcommand{\SelfInflicted}{184}
\newcommand{\SelfShare}{41}
\newcommand{\DefWordSearched}{394}
\newcommand{\DefWordThreat}{271}
\newcommand{\DefWordThreatShare}{69}
\newcommand{\DefWordOfThreats}{52}
\newcommand{\DefWordFalse}{123}
\newcommand{\PhaseThreat}{194}
\newcommand{\PhaseThreatOfThreats}{37}
\newcommand{\PhaseGateRecall}{44}
\newcommand{\PhaseGateMissed}{109}
\newcommand{\GatePrecision}{63}
\newcommand{\GateRecall}{77}
\newcommand{\GateMissed}{119}
\newcommand{\GateFalse}{239}
\newcommand{\GatePawnOnly}{176}
\newcommand{\GatePawnOnlyShare}{74}
\newcommand{\DefWordPawnOnly}{89}
\newcommand{\PassThreat}{34}
\newcommand{\PassNot}{14}
\newcommand{\Quiet}{316}
\newcommand{\QuietQuiz}{73}
\newcommand{\QuietQuizShare}{23}
\newcommand{\QuietBoard}{%
\chessboard[setfen={r2qr1k1/1p2bp1p/2p3pB/p2n1p2/1n1P4/1Q2PN2/PP1N1PPP/2R2RK1 w - - 2 16}, boardfontsize=14pt, pgfstyle=straightmove, arrow=to, linewidth=0.07em, color=red!75!black, markmoves={f3-e5}, pgfstyle=straightmove, arrow=to, linewidth=0.07em, color=green!50!black, markmoves={a2-a3}]
}
\newcommand{\QuietFen}{\fen{r2qr1k1/1p2bp1p/2p3pB/p2n1p2/1n1P4/1Q2PN2/PP1N1PPP/2R2RK1 w - - 2 16}}
\newcommand{\QuietPlayed}{16.~\symknight{}e5}
\newcommand{\QuietDrop}{16}
\newcommand{\QuietLine}{16.~a3 a4 17.~\symqueen{}d1 \symknight{}d3 18.~\symrook{}b1 \symqueen{}b6 19.~b3 \symknight{}c3}
\newcommand{\QuietUci}{\uci{a2a3 a5a4 b3d1 b4d3 c1b1 d8b6 b2b3 d5c3}}
\newcommand{\QuietTheme}{\texttt{middlegame}}
\newcommand{\QuietOptions}{16.~a3 (0.7), 16.~\symrook{}a1 ($-$0.1), 16.~a4 ($-$0.4)}


\title{Do the decks built from his games ask what they say they ask?\\[0.3em]
\large The misses deck's theme words and the guess-the-plan quiz, read against their own engine lines}
\author{Chess Coach research lab}
\date{}

\begin{document}
\maketitle

\begin{abstract}
Two practice decks are built from the player's own games. The misses deck files each position
where he went wrong under a theme word, which orders a session and, for most words, names the
skill a solve counts toward. Reading each of its \Scored{} engine-checked positions by what its
own stored lines show, \DefensiveShare\% are clear defensive misses (he had to stop something)
and \AttackingShare\% clear missed wins. The deck files \PhaseShare\% of its positions under a
phase name, among them \PhaseDefensive{} clear defensive misses and \PhaseAttacking{} clear
missed wins, and only \DefRecall\% of the clear defensive misses carry its defensive-move word.
The guess-the-plan quiz has \Quizzes{} positions; by the app's own rule \TwoRight{} of them have
two right answers out of three, the best option is a capture, a check or a two-pawn lead in
\ForcingOrMaterial, and \PlanQuestions{} ask a plan question with one right answer.
\end{abstract}

\section{The question}

The misses deck is the player's own mistakes as practice. Each position carries a theme word
from the app's labeller of the best move (\texttt{themeOfBest}): a motif when one of its detectors
fires (a mate in one, a back-rank mate, a fork, a raid on f7, a free piece, a defensive move),
else the phase. The word interleaves a session, and for the words the learner knows it is the
skill a solve or a miss is counted toward. The guess-the-plan quiz shows a position from his
games with three engine moves and asks which is the plan; an option within \Margin{} centipawns
of the best counts as right. The detector-themes project scored the labeller on puzzles, which
are tactics by construction. Here it is read on the deck itself (app commit
\texttt{\AppCommit}), against the deck's own engine lines.

\section{Method}

The rules were written in the project's README before any of these numbers was computed. What
the best move did is read from its stored line (the routine's deep pass, five lines at four
million nodes) by the board code alone: it \emph{mates} (the line ends in mate or scores a
mate), \emph{wins material} (his balance, counting 1, 3, 3, 5, 9, is at least two points up at
the line's end, the lower of the last two balances when it ends on a capture), or neither. What
his move cost is read the same way from its own stored line, or from the move and the deck's
punishing line: it \emph{allowed mate}, \emph{lost material} (two points or more), or neither.
A miss is \emph{defensive} when the best move wins nothing and his move allowed mate or lost
material, \emph{attacking} when the best move mates or wins and his move lost nothing, and
\emph{mixed} otherwise: \Both{} where both happened and \Neither{} (\NeitherShare\%) where the
stored lines show no material change either way, so material cannot judge them. The
\Deviations{} opening deviations, which have no engine lines, are left out.

\section{The misses deck}

\begin{table}[h]
\centering
\small
\ThemeTable
\caption{The deck's theme words: the learner skill each feeds, how its positions read by the
stored lines, and whether the word clearly fits by the README's rule (defensive move: a
defensive miss; hanging piece and fork: the best move wins or mates; mate words: it mates).
Phase words name no motif and are not scored.}
\end{table}

\paragraph{Most positions name no motif.} \PhaseItems{} of the \Items{} positions
(\PhaseShare\%) are filed under a phase, \DefItems{} (\DefShare\%) as a defensive move, and
\MotifItems{} (\MotifShare\%) under a tactic. The phase-word positions are not all formless:
\PhaseDefensive{} are clear defensive misses and \PhaseAttacking{} clear missed wins. On those
missed wins, the move describer the explanations use names nothing for \PhaseAttackingNothing\%.

\paragraph{The defensive-move skill sees a minority of his defensive misses.} Of the
\Defensive{} clear defensive misses, \DefFitK{} (\DefRecall\%) carry the defensive-move word and
\DefensiveElsewhere{} carry another, mostly a phase. Of the \DefItems{} positions the word is
on, \DefFitK{} (\DefFit\%) are clear defensive misses and \DefAttacking{} read as missed wins;
the rest are mixed. The labeller calls a move defensive when one of his pieces was loose and the
move protected it, which misses the threats that are not a loose piece, such as a mate.

\paragraph{Endgame positions count toward conversion.} The learner counts every position filed
under \texttt{endgame} toward \EndgameCompetency: \EndgameItems{} positions, \EndgameDefensive{}
of them clear defensive misses, where he was not converting anything.

\paragraph{The motif words.} A mate in one mates in \MateFit{} positions and a fork wins or mates
in \ForkFit. The free piece wins two points or more in \HangFit\% of its \HangItems{} positions;
\HangPawn{} more win exactly a pawn, which the two-point rule does not count, so with them it is
\HangFitPawn\%.

\begin{figure}[h]
\centering
\begin{minipage}[t]{0.47\textwidth}\centering
\DefBoard\\[0.3em]
\footnotesize Filed under \DefTheme, feeding \DefSkill
\end{minipage}\hfill
\begin{minipage}[t]{0.47\textwidth}\centering
\WinBoard\\[0.3em]
\footnotesize Filed under \WinTheme, feeding \WinSkill
\end{minipage}
\caption{Two positions from his games, drawn from White's side; he is the side to move. Red is
his move, green the engine's.}
\end{figure}

\paragraph{A mate threat filed under a phase.} In \DefFen{} he played \DefPlayed, which gave
away \DefDrop{} points of winning chances to \DefPunish. The engine's move was \DefLine{}
(\DefUci). The deck files the position under \DefTheme.

\paragraph{A win filed under conversion.} In \WinFen{} he played \WinPlayed{} (\WinDrop{}
points), where \WinLine{} (\WinUci) wins material. The deck files it under \WinTheme, so a solve
counts toward \WinSkill.

\section{The guess-the-plan quiz}

\begin{figure}[h]
\centering
\QuizBoard
\caption{A quiz position: the options are \QuizOptions, in pawns from his side. The best, a
recapture, leads the next by \QuizLead{} pawns.}
\end{figure}

By the app's rule, \OneRight{} quizzes have one right answer and \TwoRight{} (\TwoRightShare\%)
have two of three, where a guess is right two times in three. The best option is a capture or a
check in \Forcing{} and leads the next option by two pawns or more in \Material: \ForcingOrMaterial{}
(\ForcingShare\%) are one or the other, a calculation question rather than a plan. In \PlayedBest{}
he had played the best option in the game. Only \PlanQuestions{} (\PlanShare\%) ask a plan question
with one right answer. In the position above (\QuizFen), the best option is \QuizLine{} (\QuizUci).

\section{What it means for the app}

\begin{enumerate}[nosep]
\item The misses deck's word should come from what the stored lines show before the detectors
are asked: a defensive miss is one, whatever the phase, and a defensive move should include
stopping a mate or a threat the board code can name, not only a loose piece.
\item A phase word should not feed a skill: the endgame positions that were defensive misses
should not count toward conversion.
\item The quiz should be drawn from positions where one option is right and the best option is
quiet, or say when a position is a calculation question.
\end{enumerate}

\section{Limits}

The rule reads material only. A capture that removes an attacker and also wins material reads as
a missed win: \DefWinsSmall{} of the \DefWins{} defensive-move positions read as wins gain
exactly two points, and some of those are such captures, where the word is right. The \Neither{}
mixed positions, whose stored lines show no material change, cannot be judged this way; a
threat search (what the opponent would win if he passed) could read them. The stored lines are
the engine's, at the depth the routine used.

\end{document}
